\section{Regular projected recourse multipliers}\label{sec:regularity}

The inverse-order term in \cref{thm:affine-recourse} pays for a private
feasibility repair using only its length. A regular optimal multiplier
projection permits a finer estimate: the mean displacement cancels to
first order under the normalized kernel. We prove an inverse-square bound
under a weighted tensor Chebyshev condition, and an
$O(r^{-(1+\beta)})$ bound under coordinatewise H\"older regularity,
including the Lipschitz endpoint $\beta=1$. Both statements concern the
unchanged private-degree-two hierarchy of \cref{def:rec-hierarchy}.
The multipliers are used in the proof; they are not variables or inputs
of the semidefinite relaxation.

Throughout this section we assume complete affine recourse and private
convexity, and retain $d$, $\Abud_b$, and $\Abud$ from
\eqref{eq:rec-budget}. Set
\begin{equation}\label{eq:reg-parameters}
 r\ge\max(w+1,d),\qquad
 m=\left\lfloor\frac{r-1}{w}\right\rfloor+1
   =\left\lceil\frac rw\right\rceil\ge2,\qquad
 \Dm=2m^2+1,
\end{equation}
and use $V_m$ from \eqref{eq:rec-Vm}. Thus
$v_b(m-1)\le r-1$ for every bag. This single degree reserve will permit
affine localizers, source shared moments, and the multiplier commutator.

\subsection{The KKT projection and its sign}\label{sec:reg-kkt}

For the stacked system \eqref{eq:rec-stacked}, write
$D_b=\widehat E_b=[E_b;0;0]$ and
\[
 F_b(u)=\min_{y\in P_b(u)}f_b(u,y).
\]
At each $u$ suppose we choose an optimal private point $y_b^*(u)$ and
nonnegative multipliers $\pi_b(u)$ satisfying
\begin{equation}\label{eq:reg-kkt}
 \begin{split}
 q_b(u)+2Q_b(u)y_b^*(u)+\widehat A_b^\top\pi_b(u)&=0,\\
 \pi_b(u)^\top\bigl(\widehat a_b+D_bu-\widehat A_by_b^*(u)\bigr)&=0.
 \end{split}
\end{equation}
Such pairs exist pointwise even for lower-dimensional or degenerate
fibers. To see this directly, let $I$ be the active rows at an optimal
point and let $g$ be its private objective gradient. If $-g$ were outside
the cone generated by the active row normals, separation from that closed
finitely generated cone would give a direction $v$ with
$\widehat A_{b,I}v\le0$ and $g^\top v<0$. For sufficiently small positive
$t$, the inactive rows remain satisfied at $y_b^*+tv$, and the active
rows are satisfied as well. Differentiability would then give a smaller
objective, a contradiction. Hence $-g=\widehat A_{b,I}^\top\pi_I$ with
$\pi_I\ge0$; extending by zero supplies \eqref{eq:reg-kkt}. Existence of
an optimizer follows from nonempty compact fibers.

Our substantive assumptions concern only the projected multipliers
\begin{equation}\label{eq:reg-projection}
 a_b^{\mathrm{dual}}(u)=D_b^\top\pi_b(u)\in\R^{B_b}.
\end{equation}
The regular projection need not arise from a bounded, continuous, or
measurable selection of the full vector $\pi_b$. No integral below uses
that full vector.

The sign in \eqref{eq:reg-projection} follows our convention
$\widehat A_by\le\widehat a_b+D_bu$. Quadratic expansion, stationarity,
and complementarity give, for every private vector $y$,
\begin{equation}\label{eq:reg-kkt-identity}
 f_b(u,y)-F_b(u)
 -\pi_b(u)^\top(\widehat a_b+D_bu-\widehat A_by)
 =(y-y_b^*(u))^\top Q_b(u)(y-y_b^*(u))\ge0.
\end{equation}
In particular this is a lower bound on $f_b(u,y)$ even if $y$ is not
feasible at $u$.

For a feasible moment family, use $h_b,\psi_b,Y_b,M_b,U_b$ from
\cref{sec:rec-conditional}, with $m$ as in \eqref{eq:reg-parameters}.
When $h_b(u)>0$, the means $\bar x_b=U_b/h_b$ and
$\bar y_b=\psi_b/h_b$ obey $\bar y_b\in P_b(\bar x_b)$.
Apply \eqref{eq:reg-kkt-identity} at $\bar y_b$ and use
\[
 \pi_b(u)^\top(\widehat a_b+D_bu-\widehat A_b\bar y_b)
 =\pi_b(u)^\top(\widehat a_b+D_b\bar x_b-\widehat A_b\bar y_b)
    +a_b^{\mathrm{dual}}(u)^\top(u-\bar x_b).
\]
The first term is nonnegative. Matrix Jensen
\eqref{eq:rec-jensen} therefore gives
\begin{equation}\label{eq:reg-conditional}
 h_b(u)F_b(u)\le\ip{H_b(u)}{M_b(u)}
   +a_b^{\mathrm{dual}}(u)^\top\bigl(U_b(u)-u h_b(u)\bigr).
\end{equation}
The correction has a plus sign and uses $U_b-u h_b$. At zero density
$M_b=U_b=0$ by \cref{lem:conditional-matrix}, so the same inequality
holds without division by zero.

\subsection{A moment bound at the required frequencies}
\label{sec:reg-half-degree}

We first give the positivity statement that will control the correction
for truncated functionals. The usual bound requiring $\abs\alpha\le r$
does not reach all its frequencies.

\begin{lemma}[Half-degree Chebyshev moments]\label{lem:half-degree}
Let $L:\Rle{x_B}{2r}\to\R$ be $\Pre_r(B)$-positive. If
$\sum_{i\in B}\lceil\alpha_i/2\rceil\le r$, then
\[
 1\pm T_\alpha\in\Pre_r(B),\qquad
 \abs{L(T_\alpha)}\le L(1).
\]
\end{lemma}

\begin{proof}
For each $k$, $1\pm T_k$ is nonnegative on the interval, and
\cref{lem:interval-sos} gives a certificate of degree at most
$2\lceil k/2\rceil$. The even and odd cases can also be read directly
from
\[
 \begin{aligned}
 1+T_{2j}&=2T_j^2,&
 1-T_{2j}&=2(1-x^2)U_{j-1}^2\quad(j\ge1),\\
 1+T_{2j+1}&=(1+x)(U_j-U_{j-1})^2,&
 1-T_{2j+1}&=(1-x)(U_j+U_{j-1})^2\quad(j\ge0),
 \end{aligned}
\]
with $U_{-1}=0$. In the odd case substitute
$1\pm x=((1\pm x)^2+(1-x^2))/2$; the resulting degree is at most
$2j+2$. For $k=0$ the certificates are constants.

Put $z_i=T_{\alpha_i}(x_i)$ and let $k=\abs B>0$. For
$\varepsilon\in\{-1,1\}$ the parity identity is
\begin{equation}\label{eq:reg-parity}
 1+\varepsilon\prod_{i\in B}z_i
 =2^{1-k}\sum_{\substack{\eta\in\{-1,1\}^B\\
                         \prod_i\eta_i=\varepsilon}}
                   \prod_{i\in B}(1+\eta_i z_i).
\end{equation}
Expansion proves the identity: averaging over the prescribed parity
retains only the empty product and the product over all coordinates.
Multiply the univariate certificates in each term. Products of their
sums of squares are sums of squares, and all generators concern distinct
coordinates, so the resulting polynomials belong to the full box
preordering. Their degree is at most
$2\sum_i\lceil\alpha_i/2\rceil\le2r$. Applying $L$ to both signs
proves the bound. For an empty bag $T_\alpha=1$ and the assertion is
immediate.
\end{proof}

For a bounded Borel function $a$ on the shared box of bag $b$, define
\[
 (\mathcal K_ba)(x)=\int K_{m,b}(x,u)a(u)\,d\arcs^{B_b}(u).
\]
The kernel has a finite polynomial expansion, so $\mathcal K_ba$ is a
polynomial. Integrating the correction in \eqref{eq:reg-conditional}
and interchanging integration with the finite-dimensional functional
gives
\begin{equation}\label{eq:reg-commutator}
 \begin{split}
 &\int a_b^{\mathrm{dual}}(u)^\top(U_b(u)-u h_b(u))\,d\arcs^{B_b}(u)
       =L_b(R_b),\\
 &R_b(x)=\sum_{i\in B_b}\left\{
 x_i\mathcal K_b(a_{bi}^{\mathrm{dual}})(x)
       -\mathcal K_b(u_i a_{bi}^{\mathrm{dual}})(x)\right\}.
 \end{split}
\end{equation}
Every tensor Chebyshev frequency in $R_b$ has individual degrees at
most $2m-2$, except possibly one coordinate with degree at most $2m-1$.
Thus
\begin{equation}\label{eq:reg-frequency-ledger}
 \abs\alpha\le2v_b(m-1)+1\le2r-1,\qquad
 \sum_i\left\lceil\frac{\alpha_i}{2}\right\rceil
      \le v_b(m-1)+1\le r.
\end{equation}
Consequently $R_b$ lies in the actual shared moment domain and every
one of its basis moments is bounded by \cref{lem:half-degree}. Membership
in the degree-$2r$ domain alone would not justify replacing $L_b(R_b)$
by a pointwise bound for a multivariate polynomial.

\subsection{The two regularity bounds}\label{sec:reg-theorems}

\begin{theorem}[Weighted tensor Chebyshev regularity]\label{thm:reg-cheb}
Suppose the pointwise KKT pairs can be chosen so that
\begin{equation}\label{eq:reg-cheb-assumption}
 a_{bi}^{\mathrm{dual}}(u)=\sum_{\alpha\in\N^{B_b}}
               a_{bi\alpha}T_\alpha(u),\qquad
 \mathcal M=\sum_{b,i,\alpha}\abs{a_{bi\alpha}}
                              \max(1,2\alpha_i)<\infty.
\end{equation}
With the parameters \eqref{eq:reg-parameters}, every feasible family
$(L_b)$ admits a law $\nu$ on the feasible set with
\begin{equation}\label{eq:reg-cheb-rounding}
 \int f\,d\nu\le\sum_bL_b(f_b)+\frac{3(\Abud+\mathcal M)}{\Dm}.
\end{equation}
Hence
\[
 0\le\fstar-\rhorec_r=\fstar-\lambda^{\mathrm{rec}}_r
 \le\frac{3(\Abud+\mathcal M)}{\Dm}
 \le\frac{3w^2(\Abud+\mathcal M)}{2r^2}.
\]
\end{theorem}

The weight in \eqref{eq:reg-cheb-assumption} is directional: for
component $i$ it weights only frequency $\alpha_i$. Because it is at
least one, the assumption implies absolute uniform convergence on the
box. Polynomial projections satisfy it; geometrically decaying tensor
coefficients do as well. We do not infer it from arbitrary multivariate
Lipschitz continuity.

\begin{theorem}[Coordinatewise H\"older regularity]\label{thm:reg-holder}
Instead suppose that for some $0<\beta\le1$ the projection satisfies
\begin{equation}\label{eq:reg-holder-assumption}
 a_{bi}^{\mathrm{dual}}(u)=\phi_{bi}(u_i),\qquad
 \norm{\phi_{bi}}_\infty\le M_{bi},\qquad
 \abs{\phi_{bi}(v)-\phi_{bi}(v')}
                      \le H_{bi}\abs{v-v'}^\beta.
\end{equation}
Then every feasible family admits a law $\nu$ on the feasible set with
\begin{equation}\label{eq:reg-holder-rounding}
 \int f\,d\nu\le\sum_bL_b(f_b)+\frac{3\Abud}{\Dm}
   +\sum_{b,i}\left\{\frac{3M_{bi}}{\Dm}
                   +H_{bi}V_m^{(1+\beta)/2}\right\}.
\end{equation}
The same error bounds both $\fstar-\rhorec_r$ and
$\fstar-\lambda^{\mathrm{rec}}_r$. In particular it is at most
\[
 \frac{3w^2}{2r^2}\left(\Abud+\sum_{b,i}M_{bi}\right)
 +\left(\frac{\sqrt3\,w}{r}\right)^{1+\beta}\sum_{b,i}H_{bi}.
\]
For fixed data the rate is $O(r^{-(1+\beta)})$, including $O(r^{-2})$
when $\beta=1$. No sharpness is asserted for $0<\beta<1$.
\end{theorem}

Coordinatewise dependence is automatic when a bag has one shared
coordinate. For larger bags it is an additional assumption on the
projection. It places no restriction on coupled polynomial master costs
$c_b$, which do not enter private stationarity. When every shared bag has
size at most one, the model separates across distinct shared coordinates;
many private blocks may still depend on the same coordinate. The scalar
case alone is therefore not a theorem for general shared networks.

\begin{proof}[Correction estimate for \cref{thm:reg-cheb}]
We extend $\kmult_j$ by zero beyond $2m-2$. For every integer $j\ge0$,
\begin{equation}\label{eq:reg-adjacent}
 \abs{\kmult_{j+1}-\kmult_j}
 \le(2j+1)(1-\kmult_1)=\frac{3(2j+1)}{\Dm}.
\end{equation}
Indeed, \eqref{eq:rec-circle-kernel} gives
$\kmult_j=\int\cos(jt)J_m(t)\,dt/(2\pi)$ for every $j\ge0$.
The cosine difference formula and
$\abs{\sin(kt)}\le k\abs{\sin t}$ for integer $k\ge1$ give
\[
 \abs{\cos(jt)-\cos((j+1)t)}
 \le2(2j+1)\sin^2(t/2)=(2j+1)(1-\cos t).
\]
The sine bound follows by telescoping $e^{2ikt}-1$ into $k$ terms of
modulus $\abs{e^{2it}-1}$. Integrating against the nonnegative normalized
circle kernel proves \eqref{eq:reg-adjacent}, including its
support-boundary indices.

For a single tensor mode $T_\alpha$ in component $i$, all coordinates
other than $i$ contribute the factor
$\prod_{j\ne i}\kmult_{\alpha_j}\le1$. Writing $k=\alpha_i\ge1$,
its one-coordinate contribution to the commutator is
\begin{equation}\label{eq:reg-mode}
 \tfrac12(\kmult_k-\kmult_{k+1})T_{k+1}
 +\tfrac12(\kmult_k-\kmult_{k-1})T_{k-1},
\end{equation}
by $xT_k=(T_{k+1}+T_{k-1})/2$. Its absolute coefficient sum is at most
$2k(1-\kmult_1)$ by \eqref{eq:reg-adjacent}. For $k=0$ the contribution
is $(1-\kmult_1)T_1$, so the bound is
$\max(1,2k)(1-\kmult_1)$ in all cases.

A mode with $k=2m-1$ can survive through the term
$\kmult_{k-1}T_{k-1}$ even though $\kmult_k=0$; it is included in
\eqref{eq:reg-mode}. Modes with $k>2m-1$ vanish, as does any mode with
an out-of-range frequency in another coordinate. Thus the commutator has
the finite support in \eqref{eq:reg-frequency-ledger}. Absolute uniform
convergence in \eqref{eq:reg-cheb-assumption} permits termwise integration;
the coefficient bound also permits summing the absolute contributions.
\Cref{lem:half-degree}, applied to the normalized shared restriction of
$L_b$, now gives
\begin{equation}\label{eq:reg-cheb-correction}
 \abs{L_b(R_b)}\le\frac3\Dm
    \sum_{i,\alpha}\abs{a_{bi\alpha}}\max(1,2\alpha_i).
\end{equation}
\end{proof}

\begin{proof}[Correction estimate for \cref{thm:reg-holder}]
Kernel normalization removes every coordinate other than $i$ in
\eqref{eq:reg-commutator}. Hence $R_b(x)=\sum_{i\in B_b}R_{bi}(x_i)$,
where the univariate polynomial
\[
 R_{bi}(x)=\int\Kj(x,u)\phi_{bi}(u)(x-u)\,d\arcs(u)
\]
has degree at most $2m-1$. Add and subtract $\phi_{bi}(x)$ in the
integrand. By \cref{lem:rec-kernel}(b),
\[
 R_{bi}(x)=\phi_{bi}(x)(1-\kmult_1)x
  +\int\Kj(x,u)(\phi_{bi}(u)-\phi_{bi}(x))(x-u)\,d\arcs(u).
\]
The first term has absolute value at most $3M_{bi}/\Dm$. The second
has absolute value at most
$H_{bi}\int\Kj(x,u)\abs{x-u}^{1+\beta}\,d\arcs(u)$.
Since $\Kj(x,u)d\arcs(u)$ is a probability law and
$t\mapsto t^{(1+\beta)/2}$ is concave for $0<\beta\le1$,
\[
 \int\Kj(x,u)\abs{x-u}^{1+\beta}\,d\arcs(u)
 \le\left(\int\Kj(x,u)(x-u)^2\,d\arcs(u)\right)^{(1+\beta)/2}
 \le V_m^{(1+\beta)/2}.
\]
This argument includes $\beta=1$, when the concavity inequality is an
equality. Put $\varepsilon_{bi}=3M_{bi}/\Dm+H_{bi}V_m^{(1+\beta)/2}$.
The polynomial $\varepsilon_{bi}-R_{bi}$ is nonnegative on the interval
and has degree at most $2m-1$. By \cref{lem:interval-sos} its certificate
has degree at most $2m$. Since $m\le r$, this certificate embeds in the
shared scalar part of (R1), giving
$L_b(R_{bi})\le\varepsilon_{bi}L_b(1)=\varepsilon_{bi}$.
This step transfers the pointwise estimate to the truncated functional;
it uses no representing measure for $L_b$. Summing over $i$ proves the
required correction bound.
\end{proof}

\begin{proof}[Completion of \cref{thm:reg-cheb,thm:reg-holder}]
We need Borel optimal private policies to realize the local values
$F_b$. Their construction is independent of the KKT pair choices.
For $\varepsilon>0$, the strictly convex function
$f_b(u,y)+\varepsilon\norm{y}_2^2$ has a unique minimizer
$y_{b,\varepsilon}(u)$ on $P_b(u)$. This minimizer is continuous in $u$.
Indeed, along any convergent parameter sequence, private compactness gives
a cluster point, and fiber Hausdorff continuity
(\cref{lem:rec-hoffman}) approximates every limit-fiber comparison point
by feasible points in the sequence of fibers. Passing to the optimality
inequalities shows that the cluster point is the unique regularized
minimizer at the limit parameter. This proves continuity.

Fix $u$ and any unregularized minimizer $z$. Regularized optimality gives
\[
 F_b(u)\le f_b(u,y_{b,\varepsilon}(u))\le
 F_b(u)+\varepsilon\norm{z}_2^2,\qquad
 \norm{y_{b,\varepsilon}(u)}_2^2\le\norm{z}_2^2.
\]
Every cluster point as $\varepsilon\downarrow0$ is therefore an optimal
point with minimal norm. The optimizer set is nonempty compact and convex,
so its minimum-norm point is unique. Thus
$y_{b,\varepsilon}(u)$ converges to that point. Taking, for example,
$\varepsilon=1/k$ proves that the limiting policy is Borel, being a
pointwise limit of continuous policies. Its cost is $F_b(u)$.

Apply \cref{lem:rec-glue} to these policies. The resulting feasible law
has objective $\sum_b\int h_bF_b\,d\arcs^{B_b}$. Integrate
\eqref{eq:reg-conditional}, use \cref{lem:rec-matrix-cost} (which requires
$d\le r$), and then use the applicable correction estimate. This proves
both rounding inequalities. Since $\fstar\le\int f\,d\nu$, taking the
infimum over feasible families gives the moment gap. \Cref{prop:rec-dual}
gives equality with the certificate supremum and certificates at every
strict level. Finally $\Dm>2r^2/w^2$ and $V_m<3w^2/r^2$ give the stated
data-dependent rate bounds.
\end{proof}

Neither theorem controls the size of the full multiplier vector. Their
constants $\mathcal M$, $M_{bi}$, and $H_{bi}$ can grow with private
dimension and active-constraint conditioning. The largest PSD block is
still $(p_b+1)\binom{r+v_b}{v_b}$, as in \eqref{eq:rec-size}; no KKT
localizers have been added. A usable numerical accuracy guarantee requires
bounds on the stated projection norms, and the theorems supply no general
algorithm to compute them.

\subsection{Sharpness and the boundary of the hypotheses}
\label{sec:reg-examples}

\begin{example}[A sharp regular affine-recourse instance]\label{ex:reg-sharp}
Let $x\in[-1,1]$ be shared, let $v,z\in[0,1]$ be private to the two bags,
and impose $z\ge x$. Minimize
\[
 f_1(x,v)+f_2(x,z)=v^2-2xv+z^2.
\]
The first private minimum is $-x_+^2$ at $v=x_+$, and the second is
$x_+^2$ at $z=x_+$, where $x_+=\max(x,0)$. Thus $\fstar=0$ and recourse
is complete. In the first bag the projection is zero, since its
constraints are fixed. In the second, multiplier $2x_+$ on
$-z\le-x$ and zero multipliers on the box rows satisfy stationarity
and complementarity, including at $x=0$ and the endpoints. The projection
is $a_2^{\mathrm{dual}}(x)=-2x_+$, which has sup norm and Lipschitz
constant two.

Use $v=(1+\xi)/2$, $z=(1+\zeta)/2$ to place private variables in the
model box. The first bag has no additional rows and the second has
$1+\zeta-2x\ge0$. This is a positive scaling of the original row, so
rescaling its multiplier leaves its projection unchanged. The only
nonconstant shared matrix coefficient is
\[
 H_{1,(1)}=\begin{pmatrix}-1&-1/2\\-1/2&0\end{pmatrix},
\]
giving $\Abud=2$. Consequently \cref{thm:reg-holder}, with $w=1$,
$d=1$, $m=r$, $M=H=2$, gives for every $r\ge2$
\[
 -\rhorec_r\le\frac{12}{2r^2+1}+2V_r
       <\frac{24}{2r^2+1}<\frac{12}{r^2}.
\]
Conversely, choose the two measures of \cref{lem:fejer-lower} for
$\varphi(x)=x_+^2$ and $n=2r$. Lift the first by $v=x_+$ and the second
by $z=x_+$, using the appropriate private affine scalings. Both lifts
are locally feasible actual measures and their separator moments agree
through degree $2r$. Their cost is
$-\int\varphi\,d\mu_++\int\varphi\,d\mu_-$, whence
\[
 -\rhorec_r\ge\frac2{27\pi(2r+2)^2}.
\]
Thus the regular affine-recourse rate is exactly $\Theta(r^{-2})$ for
this fixed rectangular hierarchy. We transfer the separator measures,
using the lift $z=x_+$ for this problem; the lift $z=(-x)_+$ from
\cref{sec:sharpness} belongs to a different problem. No inclusion between
the rectangular and total-degree cones is used.
\end{example}

\begin{example}[Strict private convexity is insufficient]
\label{ex:strict-convex-insufficient}
For $x\in[-1,1]$, consider the scalar private problem
\[
 \min_{\abs x\le z\le1}(z^2+z)=x^2+\abs x.
\]
It has complete affine recourse, a constant positive-definite private
Hessian, and the Lipschitz optimizer $z^*(x)=\abs x$. For $0<x<1$,
the only active non-box row is $-z\le-x$, whose multiplier is $2x+1$;
the projection is $-(2x+1)$. For $-1<x<0$, the active row is $-z\le x$,
whose multiplier is $1-2x$, giving projection $1-2x$. The box rows are
inactive on both punctured neighborhoods of zero. Hence the projection
has right limit $-1$ and left limit $1$. No choice of multipliers at zero
can make it continuous. Strict private convexity and a Lipschitz primal
policy therefore do not imply either regularity assumption of this
section. This example tests the assumptions; \cref{thm:affine-sharp}
separately proves an actual inverse-order sparse gap under the general
affine-recourse hypotheses.
\end{example}

\subsection{Comparison with parametric quadratic programming}
\label{sec:reg-parametric}

Classical multiparametric QP results give piecewise-affine primal and
multiplier policies under appropriate nondegeneracy assumptions
\cite{tondel2003-an-algorithm-for-multi-parametric,baotic2016-gradient-value-function}.
They supply one useful route to our regularity assumptions. We record
its elementary argument to make the closed-box boundary precise.

\begin{proposition}[A sufficient scalar QP regime]\label{prop:reg-qp}
For each bag, suppose $Q_b$ is a constant positive-definite matrix,
$q_b$ is affine, and the affine fibers remain nonempty on an open
neighborhood of the closed shared box. Suppose that, at the unique
private optimizer throughout this neighborhood, all active rows of
$\widehat A_b$, including active private box rows, are linearly independent
(LICQ). Then the primal optimizer and the full KKT multiplier are
continuous piecewise-affine functions on the closed box and are
Lipschitz there. In particular, if each nonempty shared bag has size one,
\cref{thm:reg-holder} applies with $\beta=1$.
\end{proposition}

\begin{proof}
Suppress the bag index. Positive definiteness gives a unique primal
optimizer. For an independent active row set $I$, the KKT system is
\[
 \begin{pmatrix}2Q&\widehat A_I^\top\\\widehat A_I&0\end{pmatrix}
 \begin{pmatrix}y\\\pi_I\end{pmatrix}
 =\begin{pmatrix}-q(u)\\\widehat a_I+D_Iu\end{pmatrix}.
\]
Its matrix is nonsingular: a homogeneous solution has
$\widehat A_Iy=0$, and multiplication of its first equation by $y^\top$
gives $2y^\top Qy=0$, hence $y=0$ and then $\pi_I=0$ by row
independence. The solution is affine in $u$. Extend $\pi_I$ by zero;
primal feasibility and $\pi_I\ge0$ describe a closed polyhedral region
on which this solution is optimal. There are only finitely many row
sets, and every parameter under the hypotheses lies in one such region.

LICQ makes the multiplier at each optimizer unique, because its
nonzero entries are supported on the active independent rows. Valid
affine policies therefore agree wherever their regions overlap. Along
any convergent parameter sequence, pass to a subsequence with one of the
finitely many affine policies. Its limit is a valid KKT pair at the
limit parameter, and uniqueness identifies it with that pair. Thus the
policies are continuous, including at region boundaries and at the
boundary of the closed box. Each line segment in the box is partitioned
into finitely many intervals by these polyhedral regions. On every
interval the policy has one affine slope. Summing along the intervals
bounds the total change by the largest slope norm times the segment
length, proving Lipschitz continuity. The projection by fixed $D^\top$
is Lipschitz as well. For scalar bags, each component is a function of
its sole coordinate, so \eqref{eq:reg-holder-assumption} follows.
\end{proof}

The neighborhood assumption avoids applying an interior sensitivity
statement to an unqualified boundary point. Strict complementarity and
a single fixed active set are not required. LICQ is a sufficient route,
rather than a hypothesis of \cref{thm:reg-cheb,thm:reg-holder}; the sharp
example above need not satisfy it everywhere. For multivariate bags,
the proposition gives a Lipschitz projection but does not establish
weighted absolute tensor Chebyshev summability or coordinatewise
dependence. Oblique active-region boundaries are therefore outside the
automatic corollary.

On each critical region the affine KKT policy eliminates the private QP
exactly, leaving $c_b(u)$ plus a quadratic function. Enumerating those
regions is an established alternative to the hierarchy; no comparison
of practical effort is proved here. Polynomial lower approximation of
recourse values is also established in the joint+marginal approach
\cite{lasserre2010-joint-marginal} and in two-stage polynomial
optimization \cite{zhong2024-two-stage-polynomial}. Sparse KKT
reformulations explicitly introduce multiplier equations and auxiliary
variables \cite{qu2024-correlatively-sparse-lagrange}. Our theorems
instead give finite-order bounds for the original hierarchy while keeping
private degree two. The quantitative link uses the KKT inequality,
the kernel commutator, and its certificate at the precise truncated
degree. It is not a new QP sensitivity theorem or a claim that smoothness
alone guarantees the stated sparse rate.
